\chapter{Quantization}

Quantization is a fundamental concept in digital signal processing
that involves mapping a continuous or high-resolution discrete signal
to a finite set of levels. It is a key step in converting analog
signals to digital form. Quantization is a non-linear operation that
introduces quantization error (quantization is a lossy process).

Quantization implies as input a discretized signal in the signal
domain (in general, \emph{analog} signal samples) and the output is a
digital signal (where the signal samples are discretized in
amplitude). However, quantization can be also used to re-quantize
already quantized signals.

In general, quantization can be considered as a injective
projection between the elements of two sets $A$ and $B$, where
$|A|>|B|$, and obviously, this is an irreversible action.

\section{The quantization error}

The quantization error is the difference between the
discrete-but-still-analog input signal and the digital output signal
\begin{equation}
  \mathbf{e}_n = \mathbf{x}_n - \hat{\mathbf{x}}_n
\end{equation}

Under a high enough number of quantization levels, quantization error
is typically modeled as an uniform random variable.

\section{The rate/distortion curve}

The ultimate goal of all quantizers is to minimize the (bit-)rate and
the quantization error, which is a tradeoff (for a given quantizer, we
cannot minimize the bit-rate and the quantization error at the same
time). In terms of Rate/Distortion (R/D), the \emph{R/D efficiency} of
a quantizer can be described plotting the distortion generated for
each (possible) number of bits/sample, that at the end, determine the
number of quantization levels. Notice that, for a given bit-rate, a
quantizer is more efficient if the distortion is smaller.

\section{Scalar quantization VS vectorial quantization}

Quantizers can process the samples individually (scalar) or in groups
(vectors). Scalar quantizers are optimal in terms of R/D only if the
samples are uncorrelated. Otherwhise, vectorial quantizers (which can
exploit such correlation) are more efficient.

\section{Uniform quantization}

An uniform quantizer use equidistant decision levels (constant
quantization step size (QSS)). The representation levels (one per
interval) are usually in the middle of each interval to minimize the
quantization error, on average.

\section{Non-unidorm quantization}

Uniform quantizers only minimize the quantization error if the input
samples follow an uniform distribution. Better 

Therefore, if we know, a priori, that
actual statistical distribution, the QSS can be different over the
dynamic range of the input signal $\mathbf{x}$.

\section{Adaptive quantization}

In adaptive quantizers the QSS is adapted depending on the 
